\chapternotnumbered{Introduction} \label{ch:Introduction}

\vspace{2ex}

Orthogonal Frequency Division Multiplexing (OFDM) is a widely adopted transmission technique
for broadband communication over frequency-selective channels. By dividing a wideband
channel into a set of narrowband orthogonal subcarriers, OFDM effectively mitigates
Inter-Symbol Interference (ISI) and enables simpler equalization in the frequency
domain.

In this project, an acoustic OFDM communication system is designed, implemented, and
experimentally validated. Acoustic channels are particularly challenging due to severe
multipath propagation, long delay spreads, Doppler effects, and hardware-induced
imperfections. These impairments make synchronization, channel estimation, and phase
tracking critical components of the receiver design.

The objective of this work is threefold. First, a complete OFDM transceiver is implemented
in Matlab, including synchronization, channel estimation, equalization, and decoding.
Second, the system performance is evaluated across a set of emulated channels exhibiting
distinct impairments such as phase rotation, delay spread, band rejection, and multi-path
fading. Finally, the system is validated in a real acoustic environment through the
transmission of an image data, using the provided microphone and speakers.

Special attention is given to practical receiver algorithms, including Least Squares (LS)
and windowed LMMSE channel estimation, pilot-based phase tracking, and concatenated
Forward Error Correction (FEC). Performance is assessed using metrics such as Bit Error
Rate (BER), spectral efficiency, and channel impulse response characteristics.